\section{The convolution obstruction}
\label{sec:obstruction}

We can now exclude every nontrivial invariant projection for the measurable
weight. The argument compares two estimates for a finite part of a
convolution: polynomial compactness prevents that part from vanishing,
whereas the ordered couplings make its complementary tail too small to
cancel it.

\subsection{A conditional convolution estimate}

We isolate the common argument so that the product construction can use
it without repeating the infinite-tail estimate.  The probability space
in the lemma may include auxiliary labels besides its endpoints.

\begin{lemma}[Conditional convolution obstruction]
\label{lem:conditional-convolution}
Let $K\ge1$ be an integer, $D\ge a>0$, $c>0$, and
$\eta,\delta\ge0$.  On a probability space with endpoints $x,y$, let
$\xi(x),\gamma(y)$ be measurable vectors in the unit ball of
$\ell^2(\Z)$, and let $F_k$ be positive finite measurable functions of
an endpoint, with $F_0=1$.  Set
\[
 A_v=\frac{\xi_v(x)}{F_v(x)},\qquad
 B_w=F_{-w}(y)\overline{\gamma_{-w}(y)},\qquad
 R(k)=\frac{F_k(y)}{F_k(x)}.
\]
Suppose, almost surely, for $0\le n\le2K$ and every $v\in\Z$,
\begin{equation}\label{eq:generic-convolution}
 \sum_{v\in\Z}A_vB_{-n-v}=0,\qquad
 0<\frac{F_{n+v}(y)}{F_v(x)}
       \le\min\{R(v),R(n+v)\}\le1.
\end{equation}
Let $E$ be an event of probability at least $c$ on which
\[
 |A_v|,|B_v|\le D\quad(|v|\le K),\qquad
 |A_0|,|B_0|\ge a.
\]
For each $|k|>K$, suppose there are events $H_k^{(1)},H_k^{(2)}$ with
\begin{equation}\label{eq:generic-tail}
 R(k)\le\eta+\mathbf1_{H_k^{(1)}}+\mathbf1_{H_k^{(2)}}\quad\text{on }E,
 \qquad
 \mathbb P(H_k^{(1)}\mid x)\le\delta,\quad
 \mathbb P(H_k^{(2)}\mid y)\le\delta.
\end{equation}
The conditional bounds are under the full probability law, before
restriction to $E$.  If $\varepsilon(K,D,a)$ is the positive constant
from Lemma~\ref{lem:polynomial-coefficient}, then
\begin{equation}\label{eq:generic-obstruction}
 c\varepsilon(K,D,a)\le(2K+1)(\eta+2\sqrt\delta).
\end{equation}
\end{lemma}

\begin{proof}
The multiplier bound in \eqref{eq:generic-convolution} and
Cauchy--Schwarz make every convolution there absolutely convergent.
Put
\[
 c_n=\sum_{\substack{|v|\le K\\|n+v|\le K}}A_vB_{-n-v}
       \quad(0\le n\le2K).
\]
On $E$, Lemma~\ref{lem:polynomial-coefficient} applied to the two
truncated Laurent polynomials gives
$\sum_{n=0}^{2K}|c_n|\ge\varepsilon(K,D,a)$.  Its expected value on
$E$ is therefore at least $c\varepsilon(K,D,a)$.

Fix $n$.  For an omitted index $v$, choose $k(v)=v$ if $|v|>K$,
and $k(v)=n+v$ otherwise.  This is a deterministic choice with
$|k(v)|>K$.  The zero convolution and the two multiplier bounds give
\begin{equation}\label{eq:masked-tail}
 \mathbf1_E|c_n|\le\eta+
 \sum_{b=1}^2\sum_{v\text{ omitted}}
 |\xi_v(x)\gamma_{n+v}(y)|\mathbf1_{H_{k(v)}^{(b)}}.
\end{equation}
The unmasked sum of products of vector coordinates is at most one,
which accounts for $\eta$.  We dropped $\mathbf1_E$ from the
nonnegative exceptional terms so that their full-law conditional
bounds remain applicable.

For the first exceptional sum, Cauchy--Schwarz on probability times
counting measure bounds its expectation by
\[
 \left(\mathbb E\sum_{v\text{ omitted}}
       |\xi_v(x)|^2\mathbf1_{H_{k(v)}^{(1)}}\right)^{1/2}
 \left(\mathbb E\sum_{v\text{ omitted}}
       |\gamma_{n+v}(y)|^2\right)^{1/2}
 \le\sqrt\delta.
\]
Indeed, condition the first factor's square on $x$, apply
\eqref{eq:generic-tail} to each deterministic $k(v)$, and use
$\sum_v|\xi_v(x)|^2\le1$.  The second factor is at most one.  For
the other exceptional sum put its indicator with
$|\gamma_{n+v}(y)|^2$ and condition on $y$, giving the same bound.
Tonelli's theorem justifies every interchange of the nonnegative sums
and integrals.  No sum of the individual exceptional probabilities is
required.  Taking expectations in \eqref{eq:masked-tail} and summing
in $n$ proves \eqref{eq:generic-obstruction}.
\end{proof}

\subsection{Application to the circle weight}

\begin{theorem}\label{thm:measurable}
The bounded operator $T_*=U\pi(f_*)$ has no projection $p\in M$ satisfying
$0<\tau(p)<1$ and $(1-p)T_*p=0$.
\end{theorem}

Suppose, for contradiction, that such a projection exists. Put
\[
 \xi(x)=p(x)e_0,\qquad \gamma(y)=(1-p(y))e_0,
 \qquad r(x)=p(x)_{0,0}.
\]
Both vector norms are at most $1$. We shall use the laws $\mathbb P_j$
of Section~\ref{sec:couplings}; their expectations are denoted by
$\mathbb E_j$.

\subsection{The exact convolution identities}

By \eqref{eq:coupling-identity}, the vector
$\operatorname{diag}_v(R_{x,y}(v))\xi(x)$ belongs to $\operatorname{ran}p(y)$.
Applying $T_*(y)^n$, which preserves that range for $n\ge0$, and pairing
with $\gamma(y)$ gives
\begin{equation}\label{eq:full-convolution}
 0=\sum_{v\in\Z}\frac{F_{n+v}(y)}{F_v(x)}
       \xi_v(x)\overline{\gamma_{n+v}(y)}
   =\sum_{v\in\Z} A_vB_{-n-v},\qquad n\ge0,
\end{equation}
where, in this section,
\[
 A_v=\frac{\xi_v(x)}{F_v(x)},\qquad
 B_v=F_{-v}(y)\overline{\gamma_{-v}(y)}.
\]
All these identities hold simultaneously almost surely under each
$\mathbb P_j$. They are absolutely convergent. Indeed, monotonicity of
$F_k$ in $k$ and coordinate order imply
\begin{equation}\label{eq:ratio-minimum}
 \frac{F_{n+v}(y)}{F_v(x)}
 \le\min\{R_{x,y}(v),R_{x,y}(n+v)\}\le1,
\end{equation}
and $\sum_v|\xi_v(x)\gamma_{n+v}(y)|\le1$ by Cauchy--Schwarz.
The comparison with $R_{x,y}(v)$ follows from
$F_{n+v}(y)\le F_v(y)$; the comparison with $R_{x,y}(n+v)$ follows
from $F_{n+v}(x)\le F_v(x)$.

\subsection{A finite convolution cannot be too small}

Fix $j$ and abbreviate $K=K_j$, $m=m_j$. Define an event $E_j$ by the
following three conditions:
\begin{align*}
 &X_i<1-3/m,\quad Y_i>3/m\quad(i\in J_j),
       \quad\text{and at least one tag in }J_j;\\
 &X_i<1-3K/N_s,\quad Y_i>3K/N_s
       \quad(i\in J_s,\ s>j);\\
 &r(x)\ge1/j,\qquad 1-r(y)\ge1/j.
\end{align*}
For $0\le n\le2K$, let
\begin{equation}\label{eq:truncated-convolution}
 c_{j,n}=\sum_{\substack{|v|\le K\\ |n+v|\le K}}A_vB_{-n-v}.
\end{equation}

\begin{lemma}\label{lem:truncation-bound}
There is a constant $c_*>0$ such that, for all sufficiently large $j$,
\begin{equation}\label{eq:finite-lower-bound}
 \mathbb P_j(E_j)\ge c_*,\qquad
 \mathbb E_j\left[\mathbf 1_{E_j}
          \sum_{n=0}^{2K_j}|c_{j,n}|\right]\ge c_*\varepsilon_j.
\end{equation}
The point after which these inequalities hold may depend on $p$.
\end{lemma}

\begin{proof}
At a current-block coordinate, a tag has probability $1/m$ and
satisfies both endpoint exclusions. Since each raw marginal density is
at most $10$, the probability of an inadmissible outcome is at most
$60/m$. Thus a nontagged admissible outcome has probability at least
$1-61/m$. Independence across the block shows that exactly one tag and
admissible nontagged outcomes at the other coordinates have probability
at least
\[
 m\frac1m(1-61/m)^{m-1}\ge e^{-90}\qquad(m\ge200).
\]
For example, the last estimate follows from
$\log(1-t)\ge-t/(1-t)$ with $t=61/m$.

The total probability excluded by the future-block inequalities is at
most
\[
 60K_j\sum_{s>j}\frac1{K_s}
 \le\frac{120K_j}{K_{j+1}}\longrightarrow0.
\]
By \eqref{eq:diagonal-interior}, $0<r<1$ Haar-almost everywhere.
Lemma~\ref{lem:digit-marginals} therefore bounds the probability lost
from the last line of the definition of $E_j$ by
\[
 H_0\mu\{r<1/j\}^{1/2}
 +H_0\mu\{1-r<1/j\}^{1/2}\longrightarrow0.
\]
Consequently $c_*=e^{-90}/2$ works for all sufficiently large $j$.

On $E_j$, rounding changes the raw endpoint inequalities by at most
$1/N_s$. Since $K\ge1$, they imply
\[
 u_i(x)<1-2K/N_s,\qquad u_i(y)>2K/N_s
       \quad(i\in J_s,\ s\ge j).
\]
For $0\le v\le K$, the bound $\alpha_i<2/N_s$ makes the floors of
$u_i(x)+v\alpha_i$ and $u_i(y)-v\alpha_i$ zero in every such block.
Earlier-block floors have absolute value at most $K+1$. Hence
\[
 1/F_v(x)\le D_j,
 \qquad F_{-v}(y)\le D_j\quad(0\le v\le K).
\]
For $v<0$ these two quantities are at most $1$, by monotonicity and
$F_0=1$. Since the vector coordinates have modulus at most $1$, all
coefficients of the Laurent polynomials
\[
 A^{[K]}(z)=\sum_{|v|\le K}A_vz^v,
 \qquad B^{[K]}(z)=\sum_{|v|\le K}B_vz^v
\]
have modulus at most $D_j$. Their constant coefficients are respectively
$r(x)$ and $1-r(y)$, both at least $1/j$ on $E_j$. The coefficient of
$z^{-n}$ in their product is $c_{j,n}$. The polynomial bound
\eqref{eq:polynomial} now gives
$\sum_{n=0}^{2K}|c_{j,n}|\ge\varepsilon_j$ on $E_j$.
Integration proves the second assertion.
\end{proof}

\subsection{Controlling every part of the infinite tail}

To contradict \eqref{eq:finite-lower-bound}, we must bound the omitted
terms of \eqref{eq:full-convolution}. The important estimate is conditional
on one endpoint. It will let us sum infinitely many terms without taking
a union bound over their exceptional events.

\begin{lemma}\label{lem:tail-attenuation}
For every $j$ and every integer $k$ with $|k|>K_j$, there are measurable
events $H_k^{(1)},H_k^{(2)}$ in the labelled coupling space such that
\begin{equation}\label{eq:conditional-bad-events}
 \mathbb P_j(H_k^{(1)}\mid x)\le e^{-L_j/400},\qquad
 \mathbb P_j(H_k^{(2)}\mid y)\le e^{-L_j/400}
 \quad\text{almost surely},
\end{equation}
and on $E_j$,
\begin{equation}\label{eq:tail-ratio-bound}
 R_{x,y}(k)\le e^{-L_j}
       +\mathbf 1_{H_k^{(1)}}+\mathbf 1_{H_k^{(2)}}.
\end{equation}
The conditional bounds hold before any restriction to $E_j$.
\end{lemma}

\begin{proof}
First suppose $K_j<|k|\le N_j/4$. A tagged coordinate in $J_j$ satisfies,
after rounding,
\[
 u_i(x)<1/m_j,\qquad u_i(y)>1-1/m_j.
\]
Here $\beta=|k|\alpha_i$ lies in $(1/m_j,1/2)$. The threshold
$1-\{k\alpha_i\}$ equals $1-\beta$ for $k>0$ and $\beta$ for $k<0$.
It is crossed in either case, so \eqref{eq:ratio-crossing} gives
$R_{x,y}(k)\le e^{-C_j}\le e^{-L_j}$ on $E_j$. Set both exceptional
events to be empty in this range.

We next cover all larger $|k|$. Put $M_s=m_s^{(4/5)m_s}$. The intervals
\[
 [N_j/4,M_j],\qquad [G_sK_s,M_s]\quad(s>j)
\]
cover $[N_j/4,\infty)$. Indeed the next left endpoint is below the
preceding right endpoint by \eqref{eq:scales}, the right endpoints tend
to infinity, and $G_sK_s\le N_s/4$ because $m_s\ge4G_s$.
Split each later interval at $N_s/4$. In its lower part the thresholds
belong to $[a_s,1-a_s]$ at every coordinate of $J_s$. In its upper part,
\eqref{eq:spread} places them in $[b_s,1-b_s]$ at at least $m_s/20$
coordinates. The initial interval uses the upper-part assertion for
$s=j$.

Thus for each fixed $k$ in this tail we can choose, deterministically,
a block $J_s$, a width $w\in\{a_s,b_s\}$, and a set $I\subset J_s$
of at least $m_s/20$ indices such that
\[
 u_i^*:=1-\{k\alpha_i\}\in[w,1-w]\quad(i\in I).
\]
The lower-part choice always has $s>j$. In either type of choice,
\begin{equation}\label{eq:crossing-intensity}
 m_s\rho_s(w)\ge L_j,
 \qquad C_s\ge L_j.
\end{equation}
For $w=b_s$ the first quantity is
$\kappa_sm_s^{1/40}\ge L_s\ge L_j$; for $w=a_s$ it is
$\kappa_s\sqrt{G_s}\ge L_{s-1}\ge L_j$.

Let $H_k$ be the event of no crossing at any index in $I$, and define
\[
 H_k^{(1)}=H_k\cap
    \left\{\#\{i\in I:u_i(x)<u_i^*\}\ge|I|/2\right\},
 \qquad H_k^{(2)}=H_k\setminus H_k^{(1)}.
\]
Conditional on $x$, the first event forbids the top label at every
index in the indicated $x$-measurable set. Conditional independence
and Lemma~\ref{lem:conditional-labels} therefore give
\[
 \mathbb P_j(H_k^{(1)}\mid x)
 \le(1-\rho_s(w)/10)^{|I|/2}
 \le e^{-m_s\rho_s(w)/400}\le e^{-L_j/400}.
\]
If the $x$-measurable set has fewer than $|I|/2$ elements, the event is
empty, so the same bound holds.

For the other conditional estimate, note explicitly that coordinate
order and the absence of crossings imply
\[
 u_i(x)<u_i^*\quad\Longleftrightarrow\quad u_i(y)<u_i^*
 \qquad(i\in I).
\]
Thus $H_k^{(2)}$ requires at least $|I|/2$ indices with
$u_i(y)\ge u_i^*$. Conditional on $y$, the bottom labels at all those
indices are forbidden. The same product estimate proves the second
bound in \eqref{eq:conditional-bad-events}. These arguments use the
original law, without conditioning on $E_j$.
Figure~\ref{fig:conditioning} displays the complementary split.
\begin{figure}[tb]
\centering
\begin{tikzpicture}[
  box/.style={draw,rounded corners,align=center,text width=5.3cm,
              minimum height=1.1cm,inner sep=6pt},
  arr/.style={-{Stealth[length=2mm]},thick}]
\node[box,text width=8.2cm] (b) at (0,0)
 {$H_k$: no listed threshold is crossed\\
  On this event, $u_i(x)<u_i^*\iff u_i(y)<u_i^*$.};
\node[box] (x) at (-3.1,-2)
 {$H_k^{(1)}$: at least $|I|/2$ of the $u_i(x)$ lie below $u_i^*$};
\node[box] (y) at (3.1,-2)
 {$H_k^{(2)}$: more than $|I|/2$ of the $u_i(y)$ lie above or at $u_i^*$};
\node[box] (cx) at (-3.1,-4)
 {Condition on all of $x$;\\
  top labels can cross.\\
  $\mathbb P_j(H_k^{(1)}\mid x)\le e^{-L_j/400}$};
\node[box] (cy) at (3.1,-4)
 {Condition on all of $y$;\\
  bottom labels can cross.\\
  $\mathbb P_j(H_k^{(2)}\mid y)\le e^{-L_j/400}$};
\draw[arr] (b.south) -- (x.north);
\draw[arr] (b.south) -- (y.north);
\draw[arr] (x) -- (cx);
\draw[arr] (y) -- (cy);
\end{tikzpicture}
\caption{The two parts of the no-crossing event are small after
conditioning on complementary endpoints. Each estimate controls the
matching masked square sum in Lemma~\ref{lem:conditional-convolution}; no
simultaneous avoidance over all indices is required.}
\label{fig:conditioning}
\end{figure}


Outside $H_k$, at least one chosen coordinate crosses. All coordinates
are ordered, so \eqref{eq:ratio-crossing} and
\eqref{eq:crossing-intensity} give $R_{x,y}(k)\le e^{-L_j}$.
On $H_k$ the ratio is at most $1$. This proves
\eqref{eq:tail-ratio-bound} and the lemma.
\end{proof}

\subsection{Completion of the contradiction}

Apply Lemma~\ref{lem:conditional-convolution} to the labelled law
$\mathbb P_j$, with
\[
 K=K_j,\quad D=D_j,\quad a=1/j,\quad c=c_* =e^{-90}/2,
 \qquad \eta=e^{-L_j},\quad \delta=e^{-L_j/400}.
\]
Equations~\eqref{eq:full-convolution} and~\eqref{eq:ratio-minimum}
verify its algebraic assumptions.  The proof of
Lemma~\ref{lem:truncation-bound} supplies the event of mass at least
$c_*$ and all central-coefficient bounds for sufficiently large $j$.
Lemma~\ref{lem:tail-attenuation} gives the tail estimate with conditional
probabilities under the full law.  Consequently
\[
 c_*\varepsilon_j
 \le(2K_j+1)\bigl(e^{-L_j}+2e^{-L_j/800}\bigr)
 \le\frac{\varepsilon_j}{j},
\]
where the last inequality is the parameter choice
\eqref{eq:polynomial}.  Since $\varepsilon_j>0$ and eventually
$1/j<c_*$, this is impossible.  This proves
Theorem~\ref{thm:measurable}.
